\documentclass[../main.tex]{subfiles}
\ifSubfilesClassLoaded{\setcounter{chapter}{2}}{}
\begin{document}

\chapter{Variabel, Tipe Data, Operator, dan Masukan/Keluaran}

\begin{subcpmk}
  \item Sub-CPMK 1.2: Mendeklarasikan identifier, mengisi nilai variabel, dan menggunakan tipe data dasar dengan benar
  \item Sub-CPMK 1.3: Menggunakan operator aritmatika, relasi, dan logika serta masukan/keluaran dasar (\texttt{input}, \texttt{print})
\end{subcpmk}

\SesiRPS{3}{1.2 dan 1.3}{$\sim$170 menit}{CPMK-1}

\section{Sarana dan prasarana}

Python 3.10+, editor, terminal; Jupyter opsional untuk latihan interaktif.

\section{Prosedur kerja}

\begin{enumerate}[leftmargin=*]
  \item Pelajari tipe data dan operator dari contoh kode.
  \item Implementasikan program kalkulator/input sederhana sesuai aktivitas.
  \item Uji dengan beberapa kasus masukan dan dokumentasikan output.
\end{enumerate}

\section{Variabel dan Tipe Dasar}

Python bersifat dinamis: tipe ditentukan oleh nilai yang di-\textit{assign}. Tipe umum: \texttt{int}, \texttt{float}, \texttt{str}, \texttt{bool}.

\begin{lstlisting}[caption=Variabel dan tipe]
a = 10
b = 3.5
nama = "Python"
aktif = True
print(type(a), type(b), type(nama), type(aktif))
# <class 'int'> <class 'float'> <class 'str'> <class 'bool'>
\end{lstlisting}

\subsection{Angka (Numbers)}

Python mendukung tiga tipe angka: \texttt{int} (bilangan bulat), \texttt{float} (bilangan desimal), dan \texttt{complex} (bilangan kompleks). Pada mode interaktif, variabel khusus \texttt{\_} menyimpan hasil ekspresi terakhir.

\begin{lstlisting}[caption=Tipe angka dan operasi dasar (Tutorial Python ss3.1.1)]
>>> 17 / 3        # pembagian float
5.666666666666667
>>> 17 // 3       # pembagian bulat ke bawah
5
>>> 17 % 3        # sisa bagi (modulo)
2
>>> 2 ** 8        # pemangkatan
256
>>> 3 + 2j        # bilangan kompleks
(3+2j)
\end{lstlisting}

\begin{konsep}
Pembagian \texttt{/} \textbf{selalu} menghasilkan \texttt{float}, bahkan untuk bilangan bulat (\texttt{4 / 2} menghasilkan \texttt{2.0}). Gunakan \texttt{//} jika hasil bilangan bulat diperlukan.
\end{konsep}

\section{String}

String adalah rangkaian karakter yang diapit oleh tanda kutip tunggal (\texttt{'...'}) atau ganda (\texttt{"..."}). String bersifat \textit{immutable}: isinya tidak dapat diubah setelah dibuat.

\begin{lstlisting}[caption=Operasi dasar string (Tutorial Python ss3.1.2)]
s = "Halo, Python!"
print(len(s))          # panjang: 13
print(s[0])            # karakter pertama: 'H'
print(s[-1])           # karakter terakhir: '!'
print(s[0:4])          # slicing: 'Halo'
print(s[6:])           # 'Python!'
print("Py" + "thon")   # penggabungan: 'Python'
print("Halo " * 3)     # pengulangan: 'Halo Halo Halo '
\end{lstlisting}

\subsection{Metode String Penting}

String memiliki banyak metode bawaan. Berikut metode yang paling sering digunakan:

\begin{lstlisting}[caption=Metode string bawaan]
teks = "  Selamat Datang di Python  "

print(teks.strip())           # hapus spasi tepi: 'Selamat Datang di Python'
print(teks.lower())           # huruf kecil semua
print(teks.upper())           # huruf besar semua
print(teks.replace("Python", "Praktikum"))  # ganti kata

kata = "apel,jeruk,mangga"
daftar = kata.split(",")      # ['apel', 'jeruk', 'mangga']
print(daftar)

print(" | ".join(daftar))     # 'apel | jeruk | mangga'
print("python".startswith("py"))  # True
print("2024".isdigit())           # True
print(teks.strip().count("a"))    # hitung kemunculan 'a'
\end{lstlisting}

\subsection{Multiline String dan Raw String}

\begin{lstlisting}[caption=String multiline dan raw]
pesan = """Baris satu
Baris dua
Baris tiga"""
print(pesan)

# Raw string: karakter \ tidak dianggap escape
path = r"C:\Users\cecep\dokumen"
print(path)  # C:\Users\cecep\dokumen
\end{lstlisting}

\section{Konversi Tipe}

Gunakan \texttt{int()}, \texttt{float()}, \texttt{str()} untuk konversi eksplisit saat menggabungkan data.

\begin{lstlisting}[caption=Konversi masukan pengguna]
teks = input("Masukkan bilangan: ")
x = int(teks)
print("Kuadrat:", x * x)
\end{lstlisting}

\begin{catatan}
\texttt{input()} \textbf{selalu} mengembalikan \texttt{str}. Konversi ke angka dilakukan secara manual dengan \texttt{int()} atau \texttt{float()}. Jika pengguna mengetik bukan angka, konversi akan menghasilkan \texttt{ValueError} --- kita pelajari cara menanganinya pada modul berkas.
\end{catatan}

\section{Operator}

\begin{itemize}[leftmargin=*]
  \item \textbf{Aritmatika:} \texttt{+}, \texttt{-}, \texttt{*}, \texttt{/}, \texttt{//}, \texttt{\%}, \texttt{**}
  \item \textbf{Relasi:} \texttt{==}, \texttt{!=}, \texttt{<}, \texttt{>}, \texttt{<=}, \texttt{>=}
  \item \textbf{Logika:} \texttt{and}, \texttt{or}, \texttt{not}
  \item \textbf{Keanggotaan:} \texttt{in}, \texttt{not in} (untuk string dan koleksi)
  \item \textbf{Identitas:} \texttt{is}, \texttt{is not} (membandingkan identitas objek, bukan nilai)
\end{itemize}

\begin{lstlisting}[caption=Operator in dan chained comparison]
nama = "Python"
print("Py" in nama)         # True
print("java" not in nama)   # True

nilai = 75
print(60 <= nilai < 80)     # True -- perbandingan berantai (chained)
\end{lstlisting}

\section{Fungsi \texttt{print} Lanjut dan \texttt{input}}

\texttt{print} dapat menerima banyak argumen dan parameter \texttt{sep} (pemisah) serta \texttt{end} (penutup baris).

\begin{lstlisting}[caption=Parameter print: sep dan end]
print("apel", "jeruk", "mangga", sep=", ")
# apel, jeruk, mangga

print("Loading", end="")
print(".", end="")
print(".", end="")
print(". Selesai")
# Loading... Selesai
\end{lstlisting}

\begin{aktivitas}
  \item Tulis program yang membaca dua bilangan bulat dan menampilkan jumlah, selisih, dan hasil kali.
  \item Baca nama dan umur (tahun), lalu cetak kalimat perkenalan satu paragraf.
  \item Eksperimen: bandingkan hasil \texttt{10 / 3} dan \texttt{10 // 3}; jelaskan perbedaannya di laporan.
  \item Gunakan metode string untuk mengubah kalimat menjadi huruf kapital semua, lalu hitung berapa kali huruf 'A' muncul.
  \item Baca sebuah kalimat, kemudian cetak setiap kata pada baris tersendiri menggunakan \texttt{split()} dan loop.
\end{aktivitas}

\begin{latihan}
  \item Jelaskan perbedaan \texttt{=} sebagai penugasan dan \texttt{==} sebagai perbandingan.
  \item Apa yang terjadi jika \texttt{input} dikonversi ke \texttt{int} tetapi pengguna mengetik huruf?
  \item Tuliskan ekspresi Python untuk menghitung luas persegi panjang dari masukan\newline \texttt{panjang} dan \texttt{lebar}.
  \item Apa perbedaan antara \texttt{s[2:5]} dan \texttt{s[2:5:2]}? Berikan contoh.
  \item Tulis ekspresi satu baris menggunakan \texttt{join} untuk menghasilkan \texttt{"a-b-c"} dari list \texttt{['a','b','c']}.
  \item Refleksi: kesalahan tipe apa yang paling sering Anda temui pada latihan ini?
\end{latihan}

\begin{asesmen}
\textbf{Kuis singkat}

\begin{enumerate}
  \item Hasil dari \texttt{5 ** 2} adalah:
  \begin{enumerate}
    \item 25
    \item 7
    \item 10
    \item 25.0 saja yang valid
  \end{enumerate}
  \item Ekspresi yang menghasilkan \texttt{True} jika \texttt{x = 15}:
  \begin{enumerate}
    \item \texttt{x < 10 and x > 20}
    \item \texttt{x >= 10 or x < 5}
    \item \texttt{not (x > 0)}
    \item \texttt{x == "15"}
  \end{enumerate}
  \item Metode yang menghapus spasi di awal dan akhir string adalah:
  \begin{enumerate}
    \item \texttt{clear()}
    \item \texttt{strip()}
    \item \texttt{split()}
    \item \texttt{remove()}
  \end{enumerate}
  \item Nilai dari \texttt{"AB" * 3} adalah:
  \begin{enumerate}
    \item \texttt{"AB3"}
    \item \texttt{"ABABAB"}
    \item \texttt{["AB","AB","AB"]}
    \item Error
  \end{enumerate}
\end{enumerate}

\textbf{Tugas terstruktur}: kumpulkan berkas \texttt{tugas02.py} yang membaca tiga nilai ujian, menghitung rata-rata, dan mencetak dengan dua angka di belakang koma (petunjuk: \texttt{round} atau format string). Tambahkan juga program yang menerima satu kalimat dari pengguna, lalu cetak jumlah kata dan jumlah karakter tanpa spasi.
\end{asesmen}

\begin{checklist}
  \item Saya dapat menyimpan nilai ke variabel dan memeriksa tipe dengan \texttt{type()}
  \item Saya dapat memakai operator aritmatika dan relasi dengan benar
  \item Saya memahami perbedaan \texttt{/} dan \texttt{//}
  \item Saya dapat membaca masukan pengguna dan mengonversi ke angka
  \item Saya dapat membuat ekspresi logika sederhana
  \item Saya dapat menggunakan \textit{slicing} pada string
  \item Saya dapat menggunakan metode string seperti \texttt{split}, \texttt{strip}, \texttt{join}, \texttt{replace}
\end{checklist}

\begin{rangkuman}
Bab ini membahas variabel, tipe dasar, operator, metode string, serta \texttt{input}/\texttt{print} sebagai fondasi program interaktif.

\textbf{Poin Kunci:} konversi tipe eksplisit; \texttt{/} vs \texttt{//}; \texttt{input} menghasilkan string; string bersifat immutable; gunakan metode string untuk manipulasi teks.

\textbf{Referensi:} {\small Python Tutorial §3.1 \url{https://docs.python.org/3/tutorial/introduction.html}}

\textbf{Kata kunci:} \code{int}, \code{float}, \code{str}, \code{bool}, \code{input}, \code{print}, \code{split}, \code{strip}, \code{join}
\end{rangkuman}

\section{Referensi sesi}

\begin{itemize}[leftmargin=*]
  \item \textbf{[6]} UIMA Silab, \textit{Modul pembelajaran pemrograman Python}.
  \item \textbf{[13]} Downey, A. B. \textit{Think Python, 2e}.
  \item \textbf{[8, 9]} \textit{The Python Tutorial}.
\end{itemize}

\end{document}
